% SC26 Artifact Description for OncoRx-Bench.
% Structure follows the official SC26 sc26repro paper-author template.
% This file intentionally contains AD material only; no AE appendix is present.
\documentclass[conference]{IEEEtran}

\usepackage{cite}
\usepackage{amsmath,amssymb,amsfonts}
\usepackage{algorithmic}
\usepackage{graphicx}
\usepackage{textcomp}
\usepackage{xcolor}
\usepackage{booktabs}
\usepackage{tagging}
\usepackage{sc26repro}

% Required for final submission: keep explanations and examples disabled.
%\usetag{explanation}
%\usetag{example}

% Replace these paper-facing placeholders with immutable public revisions when
% the submission policy and final release revisions are fixed.
\newcommand{\CodeRevision}{\texttt{[FINAL CODE URL / REVISION]}}
\newcommand{\DatasetRevision}{\texttt{[FINAL DATASET URL / REVISION]}}
\newcommand{\DatasetDOI}{\texttt{[FINAL ARCHIVAL DOI, IF ASSIGNED]}}
\newcommand{\UOTDCommitA}{\texttt{e4ba3722b5505cfa587b}}
\newcommand{\UOTDCommitB}{\texttt{30032b8896d86baf8092}}
\newcommand{\ReleaseHashA}{\texttt{97a6c2a22eacf0d68e1762051fbe57b9}}
\newcommand{\ReleaseHashB}{\texttt{43f24c032a014eadd64f7e36aecf952f}}

\begin{document}

\twocolumn[%
{\begin{center}
\Huge
Appendix: Artifact Description
\end{center}}
]

%%%%%%%%%%%%%%%%%%%%%%%%%%%%%%%%%%%%%%%%%%%%%%%%%%%%%
% AD Appendix only
%%%%%%%%%%%%%%%%%%%%%%%%%%%%%%%%%%%%%%%%%%%%%%%%%%%%%

\appendixAD

\section{Overview of Contributions and Artifacts}

\subsection{Paper's Main Contributions}

\begin{description}
\item[$C_1$] A deterministic adapter that materializes generator-facing drug,
regimen, and condition views from a pinned Unified Oncology Treatment Database
(UOTD) release, with source hashes, provenance, and a fail-closed regimen
projection audit.

\item[$C_2$] A five-category, 20-subcategory oncology-medication extraction
taxonomy with fixed quotas and 347 versioned clinical-text templates.

\item[$C_3$] A standard-library-only generation and validation pipeline that
produces a deterministic 2,000-row benchmark, exact evidence offsets, grounded
annotations, and a template-disjoint 1,600/400 train/test split.

\item[$C_4$] Dataset interfaces for medication recognition and normalization,
prescribing-attribute extraction, contextual interpretation, and regimen
resolution, accompanied by category-level profiles and release manifests.
\end{description}

\subsection{Computational Artifacts}

\begin{description}
\item[$A_1$] \textbf{OncoRx-Bench reproducibility bundle.} Generator,
configuration, schemas, UOTD adapter, checked-in knowledge views, strict
validator, split/profile utilities, regression tests, and manifests.
Repository/revision: \CodeRevision. No archival DOI is assigned in this draft.

\item[$A_2$] \textbf{OncoRx-Bench dataset release.} Canonical and split JSONL
and CSV files, validation report, dataset profile, template assignments, and
release/split manifests. Dataset revision: \DatasetRevision. Archival DOI:
\DatasetDOI.
\end{description}

For dual-anonymous paper review, identity-bearing repository locations are
represented by revision placeholders. The final AD should replace each
placeholder with an immutable revision. The upstream UOTD input is pinned at
commit \UOTDCommitA\allowbreak\UOTDCommitB.

\begin{center}
\footnotesize
\begin{tabular}{@{}p{0.08\columnwidth}p{0.20\columnwidth}p{0.58\columnwidth}@{}}
\toprule
\textbf{ID} & \textbf{Supports} & \textbf{Related paper elements} \\
\midrule
$A_1$ & $C_1$--$C_3$ & Knowledge-grounded generation pipeline figure;
source-role, contracts, taxonomy, controls, and software-module tables;
integrity results. \\
\addlinespace
$A_2$ & $C_2$--$C_4$ & Dataset-composition, category-profile,
representative-instance, and benchmark-interface tables. \\
\bottomrule
\end{tabular}
\end{center}

%%%%%%%%%%%%%%%%%%%%%%%%%%%%%%%%%%%%%%%%%%%%%%%%%%%%%%%%
\section{Artifact Identification}
%%%%%%%%%%%%%%%%%%%%%%%%%%%%%%%%%%%%%%%%%%%%%%%%%%%%%%%%

\newartifact

\artrel

Artifact $A_1$ operationalizes $C_1$--$C_3$: its UOTD adapter establishes
knowledge lineage, its configuration and templates instantiate the taxonomy,
and its release driver connects generation, validation, splitting, profiling,
and byte comparison.  It supports reproducibility claims, not clinical
correctness.

\artexp

A successful check regenerates 2,000 records byte-for-byte.  The knowledge
views contain 241 drug rows, 492 regimen rows, 804 condition associations, and
2,151 projection-audit decisions.  Validation reports zero schema, quota,
identifier, grounding, offset, evidence, signature, placeholder, or duplicate
errors.  The 1,600/400 split has no shared source template; the canonical
JSONL SHA-256 is \ReleaseHashA\allowbreak\ReleaseHashB; and all 24 tests pass.
These outcomes establish artifact identity and structural integrity, not
clinical adjudication, real-note generalization, or treatment safety.

\arttime

Budget 2--5 minutes for setup, less than 1 minute for release checking, and
less than 1 minute for tests and manifest review.  On the verification host,
these steps took 0.93 and 3.14 seconds, respectively.  A network-dependent UOTD
view rebuild may require another 1--5 minutes.

\artin

\artinpart{Hardware}

A 64-bit CPU, 1~GiB available memory, and 100~MiB temporary storage suffice.
No accelerator, cluster, interconnect, or privileged service is required;
network access is used only by the optional UOTD download check.

\artinpart{Software}

The frozen environment is CPython 3.9.6 and $A_1$ at \CodeRevision.  All
runtime paths use only the standard library; \texttt{requirements.txt}
declares no external package.  A POSIX-like shell is optional.

\artinpart{Datasets / Inputs}

The normative inputs are \texttt{drug\_table.csv},
\texttt{regimen\_table.csv}, and \texttt{Conditions\_And\_Regimens.csv} under
\texttt{data/knowledge}.  Provenance, hashes, and projection decisions are
checked in.  The optional adapter check uses UOTD commit
\UOTDCommitA\allowbreak\UOTDCommitB.  The full synonym inventory is not
republished, and the derived views do not supersede upstream terms.

\artinpart{Installation and Deployment}

Clone or unpack $A_1$ and select CPython 3.9.6.  No compilation, container,
database, credentials, or service is needed.  The repository currently has no
source-code license; this must be resolved or retained explicitly as an
archival reuse limitation.

\artcomp

The workflow and dependencies are:

\begin{description}
\item[$T_1$] Verify the checked-in knowledge-view and provenance hashes.
\item[$T_2$] Using random seed 42, instantiate the 20 configured subcategories
and 347 templates to produce exactly 2,000 JSONL records.
\item[$T_3$] Run strict schema, grounding, evidence, quota, and duplicate
validation; any error terminates the workflow with a nonzero status.
\item[$T_4$] Construct template-grouped 80/20 splits and recompute exports,
profiles, LaTeX statistics, and manifests.
\item[$T_5$] Byte-compare every regenerated output with the archive.
\end{description}

Thus, $T_1\rightarrow T_2\rightarrow T_3\rightarrow T_4\rightarrow
T_5$.  The single entry point is
\texttt{python3 scripts/reproduce\_release.py --check}; the accompanying test
entry point is \texttt{python3 -m unittest discover -s tests -v}.  Because the
pipeline is deterministic, no statistical repetitions are used; tests also
check cross-process byte identity under different \texttt{PYTHONHASHSEED}s.

\artout

Inspect the validation, split, profile, and release JSON manifests.  A pass has
zero errors, exact quotas, zero template overlap, and matching SHA-256 values.
The profile generates \texttt{paper/generated\_stats.tex}, tying paper tables
to the canonical JSONL.

\newartifact

\artrel

Artifact $A_2$ supports $C_2$ and $C_4$ and is the output of $C_3$.  Its
annotations, partitions, and profile underlie the paper's dataset tables and
examples.  No LLM experiment is reported, so it reproduces the resource and
descriptive statistics, not model results.

% Manual column break keeps the final two-column page visually balanced.
\newpage
\artexp

Expect 2,000 synthetic rows in five categories and 20 subcategories, 3,525
drug objects, 1,286 populated prescribing-information objects, 444 regimen
objects, 170 unique drugs, 176 unique regimens, and 347 templates.  Full/train/
test counts are 2,000/1,600/400, validation has zero errors, and the canonical
hash is \ReleaseHashA\allowbreak\ReleaseHashB.

\arttime

Budget less than 2 minutes to obtain the 6.5~MiB bundle and 1--2 minutes to
check manifests, profiles, and examples.  No model execution is involved.

\artin

\artinpart{Hardware}

No specialized hardware is required beyond storage for a 6.5~MiB text bundle.

\artinpart{Software}

Any JSONL/CSV reader suffices.  CPython 3.9.6 and $A_1$ are needed only for
validation or regeneration; no Hugging Face package is required.

\artinpart{Datasets / Inputs}

$A_2$ at \DatasetRevision contains canonical and split JSONL/CSV plus manifests
and profiles.  Its text is synthetic and contains no patient records.  It is
not clinically adjudicated or intended for patient care; redistribution
remains subject to the notices and upstream terms in $A_1$.

\artinpart{Installation and Deployment}

Download the immutable $A_2$ revision; static inspection needs no deployment.
For regeneration, use the matching $A_1$ revision recorded in the release
manifest.

\artcomp

The dataset-level workflow is $D_1\rightarrow D_2\rightarrow D_3$:

\begin{description}
\item[$D_1$] Verify the downloaded files against the SHA-256 values in
\texttt{release\_manifest.json}.
\item[$D_2$] Confirm row counts, exhaustive split IDs, and zero template
overlap.
\item[$D_3$] Compare the profile and paper-example IDs, or invoke $A_1$ for
strict validation and regeneration.
\end{description}

Fixed parameters are the 2,000-row quotas, seed 42, nested schema, and
template-grouped subcategory-stratified split.  No repetitions are applicable.

\artout

Compare manifests and the profile with the expected values and paper tables.
Checksums establish identity; validation and split checks establish structural
consistency, not clinical validity or transfer to real notes.

\end{document}
